\documentclass{article}
\usepackage{amsmath, amssymb, amsthm}
\usepackage{hyperref}
\usepackage{geometry}
\geometry{margin=1in}

\title{Erd\H{o}s Problem \#291}
\date{Last updated: 2025-08-31}
\author{Source: erdosproblems.com}

\begin{document}
\maketitle

\noindent\textbf{Status:} OPEN \quad \textbf{No prize}

\noindent\textbf{Tags:} number theory, unit fractions

\medskip
\noindent\textbf{Problem Statement:}

Let $n\geq 1$ and define $L_n$ to be the least common multiple of $\{1,\ldots,n\}$ and $a_n$ by\[\sum_{1\leq k\leq n}\frac{1}{k}=\frac{a_n}{L_n}.\]Is it true that $(a_n,L_n)=1$ and $(a_n,L_n)>1$ both occur for infinitely many $n$?

\bigskip
\noindent\textbf{Remarks:}

Steinerberger has observed that the answer to the second question is trivially yes: for example, any $n$ which begins with a $2$ in base $3$ has $3\mid (a_n,L_n)$.

More generally, if the leading digit of $n$ in base $p$ is $p-1$ then $p\mid (a_n,L_n)$. There is in fact a necessary and sufficient condition: a prime $p\leq n$ divides $(a_n,L_n)$ if and only if $p$ divides the numerator of $1+\cdots+\frac{1}{k}$, where $k$ is the leading digit of $n$ in base $p$. This can be seen by writing\[a_n = \frac{L_n}{1}+\cdots+\frac{L_n}{n}\]and observing that the right-hand side is congruent to $1+\cdots+1/k$ modulo $p$. (The previous claim about $p-1$ follows immediately from Wolstenholme's theorem.)

This leads to a heuristic prediction (see for example a preprint of Shiu \cite{Sh16}) of $\asymp\frac{x}{\log x}$ for the number of $n\in [1,x]$ such that $(a_n,L_n)=1$. In particular, there should be infinitely many $n$, but the set of such $n$ should have density zero. Unfortunately this heuristic is difficult to turn into a proof.

Wu and Yan \cite{WuYa22} have proved, conditional on $\frac{1}{\log p}$ being linearly independent over $\mathbb{Q}$ for any finite collection of primes $p$ (itself a consequence of Schanuel's conjecture), that the set of $n$ for which $(a_n,L_n)>1$ has upper density $1$.

\bigskip
\noindent\textbf{References:}

[Sh16] P. Shiu, The denominators of harmonic numbers. arXiv:1607.02863 (2016).
 
 [WuYa22] Wu, Bing-Ling and Yan, Xiao-Hui, On the denominators of harmonic numbers. {IV}. C. R. Math. Acad. Sci. Paris (2022), 53--57.

\bigskip
\noindent\small{Source: \url{https://www.erdosproblems.com/291}}
\end{document}
